\documentstyle[%


righttag%


,11pt%


,amssymb%


,russian%               remove in Eng.


,izvru%                 Set izven% in Eng.


]{amsart}





%% Size of brackets


%% commas separation; change for point separation in English





\newcommand{\vraisup}{\operatornamewithlimits{vrai~sup}}


\newcommand{\df}{\overset{\mathrm{def}}{=}}


\newcommand{\tts}[1]{}





\theoremstyle{plain}


\theorembodyfont{\it}


\newtheorem{lemnonum}{\hskip\parindent Лемма}


\renewcommand{\thelemnonum}{}





\theoremstyle{definition}


\theorembodyfont{\rm}


\newtheorem{examnonum}{\hskip\parindent Пример}


\newtheorem{notenonum}{\hskip\parindent Замечание}


\newtheorem{assrtnonum}{\hskip\parindent Утверждение}


\renewcommand{\theexamnonum}{}


\renewcommand{\thenotenonum}{}


\renewcommand{\theassrtnonum}{}





%\hoffset=-15mm%                  For Eng.


 


\setcounter{page}{79}


\god{1999}


\nomer{4~(443)} %                        Remove in Eng.


%\nomer{Vol.\,43, No.\,4}%               Set in Eng.


%\str{\pageref{begin}--\pageref{end}}%  Set in Eng.


\udk{517.968}





\author{\it Ю.Н.\,Смолин}


% NB:   ^^ remove in Eng.


\title{Об одном методе получения экспоненциальной \\


оценки решения уравнения Вольтерра}





\date{}


%\pagestyle{myheadings}%         Set in Eng.


\pagestyle{plain} %              Remove in Eng.


\begin{document}


%\thispagestyle{plain}%          Set in Eng.





\maketitle


\label{begin}





%% NB Adjust the size of brackets





Предлагаемый метод состоит в последовательном прохождении следующих этапов.





1. Матричное ядро уравнения


\begin{equation}


x(t)=\int_0^t K(t,s)x(s)ds+f(t),\quad t\in[0,\infty[,


\end{equation}


имеющее экспоненциальный порядок роста, заменяется матрицей


\begin{equation}


\overline K(t,s)\df K(t,s)\exp(\beta(t-s)),


\end{equation}


порядок роста которой больше некоторого указанного ниже положительного числа.





2. Строится матрица $Q(\cdot,\cdot)$ --- главная часть


ядра $\overline K(\cdot,\cdot)$ 


так, что экспоненциальная оценка ее резольвенты может быть найдена с


нужной степенью точности, а получающаяся при этом добавка


$$


H(t,s)\df \overline K(t,s)-Q(t,s)


$$


не оказывает на нее существенного влияния. Tаким образом, оказывается


известной оценка резольвенты ядра $\overline K(\cdot,\cdot)$.





3. С использованием (2) определяется экспоненциальная оценка резольвенты


ядра $K(\cdot,\cdot)$.





4. Искомая оценка находится с помощью формулы решения уравнения (1)


\begin{equation}


x(t)=\int_0^t R^K(t,s)f(s)ds+f(t),\quad t\in[0,\infty[.


\end{equation}





Описанный метод в своей существенной части опирается на введенное


В.Р.\,Винокуровым [1] ассоциированное произведение матриц. Он может


быть применен, в частности, к


уравнению (1) с почти-периодическим по двум аргументам ядром; в этом


случае строится ``близкая'' к $\overline K(\cdot,\cdot)$


периодическая матрица $Q(\cdot,\cdot)$, оценка резольвенты которой


может быть найдена с любой степенью точности [2].





Приведенная в данной заметке теорема 2 дает возможность построить 


нужную $Q(\cdot,\cdot)$ с первой попытки, избавляя исследователя от


необходимости раз за разом строить эту матрицу,


находить экспоненциальную оценку ее резольвенты и уже с ее помощью


оценивать малость добавки. Это особенно обременительно, когда получение


экспоненциальной оценки резольвенты ядра $Q(\cdot,\cdot)$ сопряжено


со значительными, пусть даже техническими, трудностями. Впрочем, если


оценку резольвенты ядра $Q(\cdot,\cdot)$ получить несложно, то


можно использовать теорему~1.





Ниже $\|\cdot\|$ --- норма в $n$-мерном векторном пространстве,


$\|A\|$ --- согласованная с ней норма $n\times n$-матрицы $A$; 


$R^K(\cdot,\cdot)$ --- резольвента ядра $K(\cdot,\cdot)$;


$\Delta\df\{(t,s):0\le s\le t<\infty\}$.


При записи соотношений между суммируемыми функциями слова ``почти всюду''


будем опускать.





Напомним [3], что $n\times n$-матрица $K(\cdot,\cdot)=(k_{ij})$,


определенная в $\Delta$, удовлетворяет локальному условию $\cal U$,


если при любом $b>0$ выполняется хотя бы одно из неравенств:\linebreak


при $1<p<\infty$


$$


\biggl(\int_0^b\biggl(\int_0^t|k_{ij}|^q ds


\biggr)^{\frac{p}{q}}dt\biggr)^{\frac{1}{p}}


<\infty\quad


\biggl(\frac{1}{p}+\frac{1}{q}=1\biggr),


$$


при $p=1$


$$


\int_0^b\vraisup\limits_{s\in[0,b]}|k_{ij}|dt<\infty.


$$





Все встречающиеся в работе ядра предполагаются удовлетворяющими локальному 


условию~$\cal U$.


При этом их резольвенты (при соответствующем $p$)  удовлетворяют тому же


условию ([4], с.\,64; [5], с.\,113);  при любой $f\in L_p^n[0,b]$ существует


(и единственно) определяемое (3) решение  $x\in L_p^n[0,b]$ уравнения (1);


$R^K(\cdot,\cdot)$ в $\Delta$ удовлетворяет уравнению


$$


R^K(t,s)=K(t,s)+\int_0^t K(t,\tau)R^K(\tau,s)d\tau


$$


и является суммой составленного из итерированных ядер ряда, сходящегося по 


норме пространства операторов, действующих в $L_p^n[0,b]$ ([6]; [7], с.420);


ядра и их резольвенты суммируемы


по совокупности переменных в области $0\le s\le t<b$ [8] (и, следовательно,


справедлива теорема Фубини). Иными словами, локальное условие $\cal U$


обеспечивает законность всех производимых выкладок.





\begin{lemnonum} Пусть при $(t,s)\in\Delta$


\begin{equation}


\overline K(t,s)=Q(t,s)+H(t,s),


\end{equation}


причем


\begin{equation}


\vraisup\limits_{(t,s)\in\Delta}\|R^Q(t,s)\|


\exp(-\alpha(t-s))\le c<\infty\quad (\alpha>0).


\end{equation}





Тогда, если


$$


\vraisup\limits_{(t,s)\in\Delta}\|H(t,s)\|<m,


$$


где


\begin{equation}


m\df\min\biggl\{\frac{\sigma}{\sigma+c},\quad \frac{\sigma\alpha}{\sigma+c}


\biggr\},


\end{equation}


$\sigma>0$ --- заданное число, то


\begin{align*}


&\vraisup\limits_{(t,s)\in\Delta}\|R^{\overline K}(t,s)\|


\exp(-(\alpha+\sigma)(t-s))<\infty, \\


&\vraisup\limits_{t\in[0,\infty[}\int_0^t\|R^{\overline K}(t,s)\|


\exp(-(\alpha+\sigma)(t-s))ds<\infty.


\end{align*}


\end{lemnonum}





\begin{thm} В условиях леммы, где ядро $\overline K(\cdot,\cdot)$


определено в {\em (2)}, справедливы оценки


\begin{align}


&\vraisup\limits_{(t,s)\in\Delta}\|R^K(t,s)\|


\exp((\beta-\alpha-\sigma)(t-s))<\infty, \\


&\vraisup\limits_{t\in[0,\infty[}\int_0^t\|R^K(t,s)\|


\exp((\beta-\alpha-\sigma)(t-s))ds<\infty.


\end{align}


\end{thm}





\begin{examnonum} Пусть в (1)


$$


K(t,s)=(-0,4\exp(3(t-s))+0,01\sin t\cos s)\exp(-4(t-s)).


$$





Возьмем $\beta=4$. Тогда


$$


\overline K(t,s)=-0,4\exp(3(t-s))+0,01\sin t\cos s.


$$


Резольвента первого слагаемого легко находится,


поэтому попробуем взять


$$


Q(t,s)=-0,4\exp(3(t-s)),\quad H(t,s)=0,01\sin t\cos s.


$$


Поскольку $R^Q(t,s)=-0,4\exp(2,6(t-s))$, то


$$


\sup\limits_{(t,s)\in\Delta}|R^Q(t,s)|\le 0,4\exp(2,6(t-s)),


$$


при этом $c=0,4$; $\alpha=2,6>1$. Положив теперь $\sigma=0,01$,


в соответствии с (6) получим $m>0,02$. Так как


$$


\sup\limits_{(t,s)\in\Delta}|H(t,s)|=0,01<m,


$$


то условия теоремы 1 выполнены. Неравенства (7), (8) соответственно дают


\begin{align}


&\sup\limits_{(t,s)\in\Delta}|R^K(t,s)|\exp(1,39(t-s))<\infty,\\


&\sup\limits_{t\in[0,\infty[}\int_0^t|R^K(t,s)|\exp(1,39(t-s))ds<\infty.


\end{align}


\end{examnonum}





Назовем нижнюю грань множества значений $\alpha$, при которых


$$


\vraisup\limits_{(t,s)\in\Delta}\|K(t,s)\|\exp(-\alpha(t-s))<\infty,


$$


(экспоненциальным) порядком роста матрицы $K(\cdot,\cdot)$.





\begin{thm} Пусть определенное в {\em (2)} ядро $\overline K(\cdot,\cdot)$


таково, что


$$


\vraisup\limits_{(t,s)\in\Delta}\|\overline K(t,s)\|


\exp(-\gamma(t-s))\le c_0<\infty,


$$


где порядок роста $\gamma\ge c_0+h+1$,


\begin{equation}


h\df\frac{-(c_0+\sigma)+\sqrt{(c_0+\sigma)^2+4\sigma}}{2},


\end{equation}


$\sigma>0$ --- заданное число. Пусть, далее, 


выполняется равенство {\em (4)}, причем


$$


\vraisup\limits_{(t,s)\in\Delta}\|H(t,s)\|<h.


$$





Тогда, если выполняется {\em (5)}, где $c\le c_0+h$,


справедливы оценки {\em (7), (8)}.


\end{thm}





\begin{notenonum} Так как при $(t,s)\in\Delta$


$$


\|Q(t,s)\|\le c_0\exp(\gamma(t-s))+h\le (c_0+h)\exp(\gamma(t-s)),


$$


то [6]


$$


\|R^Q(t,s)\|\le (c_0+h)\exp((\gamma+c_0+h)(t-s)).


$$


Сравнивая это неравенство с (5), видим, что условие


$c\le c_0+h$ естественно.


\end{notenonum}





Проиллюстрируем теорему 2, вернувшись к рассмотренному примеру. Имеем


$$


\sup\limits_{(t,s)\in\Delta}|K(t,s)|\le 0,41\exp(-(t-s)),


$$


и, следовательно, $c_0=0,41$.





Примем $\sigma=0,01$. Тогда в соответствии с (11)


$0,02<h<0,03$. Поскольку должно быть $\gamma\ge c_0+h+1$,


то берем $\beta=3$. Получим


$$


\overline K(t,s)=-0,4\exp(2(t-s))+0,01\sin t\cos s\exp(-(t-s)).


$$





Tак как


$$


\sup\limits_{(t,s)\in\Delta}|0,01\sin t\cos s\exp(-(t-s))|\le 0,01<h,


$$


то положим


$$


Q(t,s)=-0,4\exp(2(t-s)),\ H(t,s)=0,01\sin t\cos s\exp(-(t-s)).


$$


Условия теоремы 2 выполнены, и остается найти оценку (5) при $c\le c_0+h$.


Имеем $R^Q(t,s)=-0,4\exp(1,6(t-s))$, и требуемое условие, очевидно,


выполняется. Поэтому справедливы неравенства (7), (8). Как нетрудно


убедиться, получаем (9), (10).





Для логической завершенности работы приведем утверждение


относительно решения уравнения (1).





\begin{assrtnonum} Если $R^K(\cdot,\cdot)$ удовлетворяет условию (8) и


$$


\vraisup\limits_{t\in[0,\infty[}\|f(t)\|\exp((\beta-\alpha-\sigma)t)<\infty,


$$


то


$$


\vraisup\limits_{t\in[0,\infty[}\|x(t)\|\exp((\beta-\alpha-\sigma)t)<\infty.


$$


\end{assrtnonum}





\begin{thebibliography}{9}





\bibitem{1}


Винокуров В.Р. {\em Аппроксимация квазилинейных интегральных уравнений 


Вольтерра алгебраическими уравнениями} //


Изв. вузов. Математика. -- 1963. -- \No\,6. -- С.\,39--48.


\bibitem{2}


Винокуров В.Р., Смолин Ю.Н. {\em Об асимптотике уравнений Вольтерра


с почти-периодическими ядрами и запаздываниями}


// ДАН СССР. -- 1971. -- Т.\,201. -- \No\,4. -- С.\,1704--1707.


\bibitem{3}


Забрейко П.П., Кошелев А.И., Красносельский М.А.


{\em Интегральные уравнения}. -- М.: Наука, 1968. -- 448~с.


\bibitem{4}


Азбелев Н.В., Максимов В.П., Рахматуллина Л.Ф. {\em Введение в теорию


функционально-дифференциальных уравнений.} -- М.: Наука, 1991. -- 280 с.


\bibitem{5}


Рахматуллина Л.Ф. {\em Линейные функционально-дифференциальные уравнения}:


Дис.\;$\dotsc$ докт. физ.-матем. наук. -- Киев, 1982. -- 280 с.


\bibitem{6}


Комленко Ю.В. {\em Необходимое и достаточное условие справедливости


теоремы об интегральных неравенствах в пространстве суммируемых функций} //


Тр. Ижевск. матем. семин. -- 1975. -- \No\,3. -- С.\,73--86.


\bibitem{7}


Канторович Л.В., Акилов Г.П. {\em Функциональный анализ.} --


2-е изд. -- М.: Наука, 1977. -- 741~с.


\bibitem{8}


Салехов Д.В. {\em Пример вполне непрерывного интегрального оператора, 


действующего из


${\cal L}_p$ в ${\cal L}_p$ с положительным ядром, не принадлежащим


${\cal L}_r$ $(r>1)$} // УМН. -- 1963. -- Т.\,18. --


Вып.\,4. -- C.\,179--182.





\end{thebibliography}





\vskip 0.5cm





\post{\it Магнитогорский педагогический}{\it Поступила}


\post{\it институт}{10.04.1998}





\label{end}


\end{document}








